% ------------------------------------------------------------------------
% file `3_juin_2023-connaitre-1-exercise-body.tex'
%   in folder `xsim/'
%
%     exercise of type `connaitre' with id `1'
%
% generated by the `connaitre' environment of the
%   `xsim' package v0.21 (2022/02/12)
% from source `3_juin_2023' on 2023/06/15 on line 70
% ------------------------------------------------------------------------
    Détermine la ou les bonnes réponses. Si tu oublies des bonnes réponses ou que tu en choisis des mauvaises, tu n'obtiens aucun point pour la question.
    \begin{enumerate}
        \item Une fonction linéaire est une fonction qui \dots
              \begin{QCM}(2)
                  \choice passe par l'origine du repère
                  \choice a une pente nulle
                  \choice est toujours croissante
                  \choice admet $0$ comme racine.
              \end{QCM}
        \item Combien de solution(s) admet l'équation $x^2+4=0$ ?
              \begin{QCM}(3)
                  \choice $2$
                  \choice $1$
                  \choice $0$
              \end{QCM}
        \item Quel est le degré du polynôme $x^6 + 5x^4 -2x^3 +23$ ?
              \begin{QCM}(3)
                  \choice $3$
                  \choice $4$
                  \choice $6$
              \end{QCM}
        \item Quelle(s) expression(s) sont équivalentes à $\dfrac{3}{\sqrt{2}}$ ?
              \begin{QCM}(4)
                  \choice $\dfrac{3\sqrt{2}}{2}$
                  \choice $\dfrac{\sqrt{6}}{2}$
                  \choice $\sqrt{\dfrac{3}{2}}$
                  \choice $\sqrt{\dfrac{9}{2}}$
              \end{QCM}
    \end{enumerate}
